\chapterbackmatter{Solutions to Weinberg Exercises}
\ExerciseEditorialNote

\begin{enumerate}
\WeinbergSolution{1}
Introduce the boundary-value functions
\[
 F^\pm(E)=\sum_\gamma\frac{|u_\gamma|^2}
 {E-E_\gamma\pm \ii0}.
\]
The Lippmann--Schwinger equation and the separable form of the interaction
give
\[
 \braket{\beta}{\alpha^\pm}
 =\delta_{\beta\alpha}
 \frac{g\,u_\beta}{E_\alpha-E_\beta\pm \ii0}
 \sum_\gamma u_\gamma^*\braket{\gamma}{\alpha^\pm}.
\]
Multiplying by \(u_\beta^*\), summing over \(\beta\), and solving the
resulting scalar equation yields
\[
 \sum_\gamma u_\gamma^*\braket{\gamma}{\alpha^\pm}
 =\frac{u_\alpha^*}{1-gF^\pm(E_\alpha)}.
\]
Thus the explicit scattering states are
\[
 \begin{split}
 \InKet{\alpha}
 &=\ket{\alpha}_0+
 \frac{g u_\alpha^*}{1-gF^+(E_\alpha)}
 \sum_\beta\frac{u_\beta\ket{\beta}_0}
 {E_\alpha-E_\beta+\ii0},\\
 \OutKet{\alpha}
 &=\ket{\alpha}_0+
 \frac{g u_\alpha^*}{1-gF^-(E_\alpha)}
 \sum_\beta\frac{u_\beta\ket{\beta}_0}
 {E_\alpha-E_\beta-\ii0}.
 \end{split}
\]
In particular,
\[
 T^+_{\beta\alpha}
 =\bra{\beta}H_{\rm int}\InKet{\alpha}
 =\frac{g\,u_\beta u_\alpha^*}
 {1-gF^+(E_\alpha)}.
\]
With the convention of Eq.~\eqref{eq:3.1.14}, the answer for the
\(S\)-matrix is therefore
\[
 \OutBra{\beta}\InKet{\alpha}
 =\delta(\beta-\alpha)
 -2\pi \ii\,\delta(E_\beta-E_\alpha)
 \frac{g\,u_\beta u_\alpha^*}
 {1-gF^+(E_\alpha)}.
\]
Since \(F^-(E)=F^+(E)^*\) for real \(E\) and real \(g\), this expression
also displays the relation between the incoming and outgoing boundary
values required by unitarity.

\WeinbergSolution{2}
At the peak of an isolated resonance, Eq.~\eqref{eq:3.8.14} gives, for
unpolarized \(e^+e^-\) beams and \(j_R=1\),
\[
 \sigma_{\rm el}^{\rm peak}
 =\frac{4\pi}{k^2}\frac{2j_R+1}
 {(2s_{e^-}+1)(2s_{e^+}+1)}B_e^2
 =\frac{3\pi}{k^2}B_e^2,
\]
where \(B_e=\Gamma(R\to e^+e^-)/\Gamma\) and
\(k=E_{\rm cm}/2=75\,{\rm GeV}\).  Using
\[
 1\ {\rm GeV}^{-2}=0.389379\times10^{-27}\ {\rm cm}^2
\]
gives
\[
 \frac{3\pi}{k^2}=1.6755\times10^{-3}\ {\rm GeV}^{-2}
 =6.523\times10^{-31}\ {\rm cm}^2.
\]
Consequently
\[
 B_e=\left(\frac{10^{-34}}{6.523\times10^{-31}}\right)^{1/2}
 =1.24\times10^{-2}.
\]
The inclusive peak cross section contains only one power of the entrance
branching fraction:
\[
 \sigma_{\rm tot}^{\rm peak}
 =\frac{3\pi}{k^2}B_e
 =8.08\times10^{-33}\ {\rm cm}^2.
\]
Thus the electronic branching ratio is about \(1.24\%\), and the total
peak cross section is about \(8.1\times10^{-33}\ {\rm cm}^2\).

\WeinbergSolution{3}
Let particle \(2\) be initially at rest, and choose the incident momentum
along \(z\):
\[
 p_1=(E_1,0,0,p),\qquad p_2=(m_2,\mathbf0),\qquad
 p_3=(E_3,q\widehat{\mathbf n}),\quad
 \widehat{\mathbf n}\cdot\widehat{\mathbf z}=\cos\theta .
\]
Momentum conservation fixes \(\mathbf p_4=\mathbf p-\mathbf q\), so
\[
 E_3=\sqrt{q^2+m_3^2},\qquad
 E_4=\sqrt{p^2+q^2-2pq\cos\theta+m_4^2}.
\]
The positive physical root \(q=q(\theta)\) is determined by
\[
 E_1+m_2=E_3+E_4.
\]
In the laboratory frame the relative speed is \(u_\alpha=p/E_1\).
After using the momentum delta function in Eq.~\eqref{eq:3.4.15}, the
remaining energy delta function has the Jacobian
\[
 \left|\frac{d(E_3+E_4)}{\dd q}\right|
 =\left|\frac q{E_3}+\frac{q-p\cos\theta}{E_4}\right|.
\]
It follows directly that
\[
 \boxed{\;
 \frac{\dd\sigma}{\dd\Omega}_{\rm lab}
 =\frac{(2\pi)^4E_1}{p}|M_{\beta\alpha}|^2
 \frac{q^2}
 {\left|q/E_3+(q-p\cos\theta)/E_4\right|}
 \;}
\]
or, equivalently,
\[
 \frac{\dd\sigma}{\dd\Omega}_{\rm lab}
 =\frac{(2\pi)^4E_1}{p}|M_{\beta\alpha}|^2
 \frac{q^2E_3E_4}
 {|q(E_3+E_4)-pE_3\cos\theta|}.
\]
The usual average over unresolved initial spins and sum over final spins
is understood.  If the final particles are identical, the result is
integrated over only one copy of phase space (or multiplied by
\(1/2!\) when the full solid angle is used).

\WeinbergSolution{4}
For every denominator in Eq.~\eqref{eq:3.5.3}, use
\[
 \frac1{E_\alpha-E_\gamma+\ii0}
 =-\ii\int_0^\infty \dd\tau\,
 e^{\ii(E_\alpha-E_\gamma)\tau},
\]
and represent the energy-conserving factor in the \(S\)-matrix by
\[
 2\pi\delta(E_\beta-E_\alpha)
 =\int_{-\infty}^{\infty}\dd t\,e^{\ii(E_\beta-E_\alpha)t}.
\]
At second order, setting \(t_1=t\) and \(t_2=t-\tau\) gives
\[
 \begin{split}
 &-2\pi \ii\delta(E_\beta-E_\alpha)
 \sum_\gamma\frac{(H_{\rm int})_{\beta\gamma}
 (H_{\rm int})_{\gamma\alpha}}
 {E_\alpha-E_\gamma+\ii0}\\
 &\qquad=(-\ii)^2\int_{-\infty}^{\infty}\dd t_1
 \int_{-\infty}^{t_1}\dd t_2\,
 \bra{\beta}H_I(t_1)H_I(t_2)\ket{\alpha} .
 \end{split}
\]
For the \(n\)th term introduce
\[
 t_1=t,\qquad
 t_{r+1}=t_r-\tau_r,\qquad \tau_r\geq0.
\]
The domain is exactly
\(t_1\geq t_2\geq\cdots\geq t_n\).  Inserting a complete set of
\(H_0\)-eigenstates between consecutive factors of \(H_I\) reproduces
all the exponentials from the energy denominators.  Hence the \(n\)th
old-fashioned term becomes
\[
 (-\ii)^n\int_{-\infty}^{\infty}\dd t_1
 \int_{-\infty}^{t_1}\dd t_2\cdots
 \int_{-\infty}^{t_{n-1}}\dd t_n\,
 \bra{\beta}H_I(t_1)\cdots H_I(t_n)\ket{\alpha} .
\]
Adding the identity term and suppressing the external matrix elements
gives precisely Eq.~\eqref{eq:3.5.8}.

\WeinbergSolution{5}
Write
\[
 G_{\rm P}(E)=\operatorname{P.V.}\frac1{E-H_0}.
\]
Both \(H_{\rm int}\) and \(G_{\rm P}(E)\) are Hermitian.  The standing-wave
transition operator satisfies
\[
 T^0=H_{\rm int}+H_{\rm int}G_{\rm P}T^0
 =(1-H_{\rm int}G_{\rm P})^{-1}H_{\rm int}.
\]
Using
\((1-AB)^{-1}A=A(1-BA)^{-1}\), one also has
\[
 T^0=H_{\rm int}(1-G_{\rm P}H_{\rm int})^{-1}=(T^0)^\dagger.
\]
Restriction to equal-energy states and multiplication by the real kernel
\(\pi\delta(E_\beta-E_\alpha)\) therefore makes
\[
 K_{\beta\alpha}
 =\pi\delta(E_\beta-E_\alpha)T^0_{\beta\alpha}
\]
Hermitian.

The identity
\[
 \frac1{x+\ii0}=\operatorname{P.V.}\frac1x-\ii\pi\delta(x)
\]
relates the outgoing resolvent to the principal-value resolvent.  If
\(X=\pi\delta(E-H_0)T^+\) on the energy shell, its integral equation is
\[
 X=K-\ii KX,\qquad X=(1+\ii K)^{-1}K.
\]
Since \(X\) is a function of \(K\), the order is immaterial, and
\[
 \boxed{\quad S=1-2\ii X
 =(1-\ii K)(1+\ii K)^{-1}.\quad}
\]
Hermiticity of \(K\) makes the Cayley transform manifestly unitary.

\WeinbergSolution{6}
For pion--nucleon scattering, couple the pion orbital angular momentum
\(\ell\) to the nucleon spin \(1/2\).  Denote
\[
 S_{\ell\pm}^{I}=e^{2\ii\delta_{\ell\pm}^{I}},\qquad
 j=\ell\pm\frac12,\qquad P=(-1)^{\ell+1},
\]
where \(I=1/2,3/2\).  The spin-nonflip and spin-flip amplitudes in a
definite isospin channel may be written
\[
 \begin{split}
 f_I(\theta)&=\frac1{2\ii k}\sum_{\ell=0}^\infty
 \left[(\ell+1)(S_{\ell+}^{I}-1)
 +\ell(S_{\ell-}^{I}-1)\right]P_\ell(\cos\theta),\\
 g_I(\theta)&=\frac1{2\ii k}\sum_{\ell=1}^\infty
 (S_{\ell+}^{I}-S_{\ell-}^{I})P_\ell^1(\cos\theta).
 \end{split}
\]
(Changing the convention for \(P_\ell^1\) reverses the irrelevant overall
sign of \(g_I\).)  Averaging over the initial proton spin and summing over
the final spin gives
\[
 \frac{\dd\sigma_I}{\dd\Omega}=|f_I|^2+|g_I|^2.
\]
The state \(\pi^+p\) has \(I_3=3/2\), and hence is pure \(I=3/2\):
\[
 \frac{\dd\sigma(\pi^+p)}{\dd\Omega}
 =|f_{3/2}|^2+|g_{3/2}|^2.
\]
Clebsch--Gordan decomposition of the elastic \(\pi^-p\) channel gives
\[
 f_{\pi^-p}=\frac{2f_{1/2}+f_{3/2}}3,\qquad
 g_{\pi^-p}=\frac{2g_{1/2}+g_{3/2}}3,
\]
and therefore
\[
 \frac{\dd\sigma(\pi^-p)}{\dd\Omega}
 =\left|\frac{2f_{1/2}+f_{3/2}}3\right|^2
 +\left|\frac{2g_{1/2}+g_{3/2}}3\right|^2.
\]
These formulas express the answer entirely through phase shifts of
definite \(j\), parity, and isospin.

\WeinbergSolution{7}
Insert a complete set of two-particle momentum--spin states between the
bra and ket and use Eq.~\eqref{eq:3.7.5} twice.  The two total-momentum
delta functions give
\(\delta^3(\mathbf p'-\mathbf p)\).  In the center-of-mass frame write
\(\mathbf p_2=-\mathbf p_1\), \(k=|\mathbf p_1|\), and
\(E=E_1(k)+E_2(k)\).  The radial energy delta function has derivative
\[
 \frac{d(E_1+E_2)}{\dd k}
 =k\left(\frac1{E_1}+\frac1{E_2}\right)
 =\frac{kE}{E_1E_2}.
\]
Thus its Jacobian converts the radial measure as
\[
 k^2\,\dd k\,
 \delta(E-E_1-E_2)\delta(E'-E_1-E_2)
 =\frac{kE_1E_2}{E}\,\delta(E'-E).
\]
This factor is cancelled exactly by the product of the two normalization
factors \((kE_1E_2/E)^{-1/2}\) in Eq.~\eqref{eq:3.7.5}.

The angular integral gives
\(\int \dd\Omega\,Y_{\ell'}^{m'*}Y_\ell^m
=\delta_{\ell'\ell}\delta_{m'm}\).  Successive use of the two
Clebsch--Gordan orthogonality relations then produces
\(\delta_{s's}\delta_{j'j}\delta_{\sigma'\sigma}\), while the species
labels give \(\delta_{n'n}\).  Altogether,
\[
 \braket{E',\mathbf p',j',\sigma',\ell',s',n'}
 {E,\mathbf0,j,\sigma,\ell,s,n}
 =\delta^3(\mathbf p')\delta(E'-E)
 \delta_{j'j}\delta_{\sigma'\sigma}
 \delta_{\ell'\ell}\delta_{s's}\delta_{n'n},
\]
which is Eq.~\eqref{eq:3.7.6}.  For identical particles the independent
integration region is only half of momentum space; the additional
\(\sqrt2\) in each state restores the same normalization.
\end{enumerate}
